% GENERATED FILE. Edit canonical lessons, then run npm run generate:books.
% canonical-source-hash: fnv1a64:b70a72ead0960973
% canonical-lessons: PT-C505-pincel, PT-C505-aguarela, PT-C505-tela, PT-C505-cavalete, PT-C505-moldura

\chapter{Pincéis e telas}
\label{ch:pt-pinceis-e-telas}
\hlchaptermodality{505}

\begin{chapteropening}
\textbf{By the end of this chapter:} \emph{I can talk about painting as a hobby.}

Complete the last lesson of chapter 505: I can talk about painting as a hobby.
\end{chapteropening}

\section[o pincel]{o pincel — paintbrush}

\label{lesson:PT-C505-pincel}



\begin{quote}
Before the new one: say the Portuguese for clay, then the Portuguese for pottery.
\end{quote}

\subsection*{You'll want to know: o pincel}
\textbf{o pincel} — \textquotedblleft{}paintbrush\textquotedblright{}.

\begin{quote}
Ela pinta com um pincel muito fino.
\end{quote}

\begin{quote}
\emph{She paints with a very fine brush.}
\end{quote}

Plural \emph{pincéis}. A hairbrush is \emph{uma escova}.

\begin{cousinweb}
\emph{pincel} comes from Catalan \emph{pinzell} \textquotedblleft{}paintbrush\textquotedblright{} (EtymDB), which goes back to Latin \emph{pēnicillum} \textquotedblleft{}a painter's brush\textquotedblright{}.
\end{cousinweb}

\begin{grammarlens}[title={What you've built}]
One more word for telling, explaining and describing at B1.
\end{grammarlens}

\subsection*{Your turn}
\begin{itemize}
\raggedright
  \item \emph{Say it:} \emph{o pincel}
  \item \emph{Say it:} \emph{o pincel}, once more
  \item \emph{Say it:} a sentence of your own with \emph{o pincel}
  \item \emph{Recall:} say the Portuguese for clay, then the Portuguese for pottery, then say \emph{o pincel} again
\end{itemize}

\begin{tcolorbox}[breakable,colback=teal!4,colframe=teal!35!black,title={Before you move on}]
What does o pincel mean? (\textquotedblleft{}Paintbrush\textquotedblright{}.) Say it once more.
\end{tcolorbox}
\section[a aguarela]{a aguarela — watercolour}

\label{lesson:PT-C505-aguarela}



\begin{quote}
Before the new one: say the Portuguese for pottery, then the Portuguese for paintbrush.
\end{quote}

\subsection*{You'll want to know: a aguarela}
\textbf{a aguarela} — \textquotedblleft{}watercolour\textquotedblright{}.

\begin{quote}
A minha tia pinta paisagens a aguarela.
\end{quote}

\begin{quote}
\emph{My aunt paints landscapes in watercolour.}
\end{quote}

In Brazil: \emph{aquarela}.

\begin{cousinweb}
\emph{aguarela} is usually traced to Italian \emph{acquerello}, built on \emph{acqua} \textquotedblleft{}water\textquotedblright{}, from Latin \emph{aqua} (Lewis \& Short).
\end{cousinweb}

\begin{grammarlens}[title={What you've built}]
One more word for telling, explaining and describing at B1.
\end{grammarlens}

\subsection*{Your turn}
\begin{itemize}
\raggedright
  \item \emph{Say it:} \emph{a aguarela}
  \item \emph{Say it:} \emph{a aguarela}, once more
  \item \emph{Say it:} a sentence of your own with \emph{a aguarela}
  \item \emph{Recall:} say the Portuguese for pottery, then the Portuguese for paintbrush, then say \emph{a aguarela} again
\end{itemize}

\begin{tcolorbox}[breakable,colback=teal!4,colframe=teal!35!black,title={Before you move on}]
What does a aguarela mean? (\textquotedblleft{}Watercolour\textquotedblright{}.) Say it once more.
\end{tcolorbox}
\section[a tela]{a tela — canvas}

\label{lesson:PT-C505-tela}



\begin{quote}
Before the new one: say the Portuguese for paintbrush, then the Portuguese for watercolour.
\end{quote}

\subsection*{You'll want to know: a tela}
\textbf{a tela} — \textquotedblleft{}canvas\textquotedblright{}.

\begin{quote}
O pintor comprou uma tela grande.
\end{quote}

\begin{quote}
\emph{The painter bought a large canvas.}
\end{quote}

In Brazil it also means a screen; in Portugal a screen is \emph{o ecrã}.

\begin{cousinweb}
\emph{tela} is inherited from Latin \emph{tēla} \textquotedblleft{}a web, cloth on a loom\textquotedblright{} (EtymDB), from \emph{texere} \textquotedblleft{}to weave\textquotedblright{}.
\end{cousinweb}

\begin{grammarlens}[title={What you've built}]
One more word for telling, explaining and describing at B1.
\end{grammarlens}

\subsection*{Your turn}
\begin{itemize}
\raggedright
  \item \emph{Say it:} \emph{a tela}
  \item \emph{Say it:} \emph{a tela}, once more
  \item \emph{Say it:} a sentence of your own with \emph{a tela}
  \item \emph{Recall:} say the Portuguese for paintbrush, then the Portuguese for watercolour, then say \emph{a tela} again
\end{itemize}

\begin{tcolorbox}[breakable,colback=teal!4,colframe=teal!35!black,title={Before you move on}]
What does a tela mean? (\textquotedblleft{}Canvas\textquotedblright{}.) Say it once more.
\end{tcolorbox}
\section[o cavalete]{o cavalete — easel}

\label{lesson:PT-C505-cavalete}



\begin{quote}
Before the new one: say the Portuguese for watercolour, then the Portuguese for canvas.
\end{quote}

\subsection*{You'll want to know: o cavalete}
\textbf{o cavalete} — \textquotedblleft{}easel\textquotedblright{}.

\begin{quote}
Pôs o cavalete perto da janela para ter luz.
\end{quote}

\begin{quote}
\emph{She put the easel near the window to get light.}
\end{quote}

Literally \textquotedblleft{}a little horse\textquotedblright{}: it stands on legs.

\begin{cousinweb}
\emph{cavalete} is usually traced to Italian \emph{cavalletto} \textquotedblleft{}a little horse, a trestle\textquotedblright{}, from Latin \emph{caballus} \textquotedblleft{}horse\textquotedblright{} (Lewis \& Short).
\end{cousinweb}

\begin{grammarlens}[title={What you've built}]
One more word for telling, explaining and describing at B1.
\end{grammarlens}

\subsection*{Your turn}
\begin{itemize}
\raggedright
  \item \emph{Say it:} \emph{o cavalete}
  \item \emph{Say it:} \emph{o cavalete}, once more
  \item \emph{Say it:} a sentence of your own with \emph{o cavalete}
  \item \emph{Recall:} say the Portuguese for watercolour, then the Portuguese for canvas, then say \emph{o cavalete} again
\end{itemize}

\begin{tcolorbox}[breakable,colback=teal!4,colframe=teal!35!black,title={Before you move on}]
What does o cavalete mean? (\textquotedblleft{}Easel\textquotedblright{}.) Say it once more.
\end{tcolorbox}
\section[a moldura]{a moldura — (picture) frame}

\label{lesson:PT-C505-moldura}



\begin{quote}
Before the new one: say the Portuguese for canvas, then the Portuguese for easel.
\end{quote}

\subsection*{You'll want to know: a moldura}
\textbf{a moldura} — \textquotedblleft{}(picture) frame\textquotedblright{}.

\begin{quote}
Comprei uma moldura para a fotografia dos avós.
\end{quote}

\begin{quote}
\emph{I bought a frame for my grandparents' photograph.}
\end{quote}

The frame of glasses is \emph{a armação}.

\begin{cousinweb}
\emph{moldura} was borrowed from Spanish \emph{moldura} (EtymDB), built on \emph{molde} \textquotedblleft{}a mould\textquotedblright{}, which goes back to Latin \emph{modulus} \textquotedblleft{}a small measure\textquotedblright{}, from \emph{modus}.
\end{cousinweb}

\begin{grammarlens}[title={What you've built}]
One more word for telling, explaining and describing at B1.
\end{grammarlens}

\subsection*{Your turn}
\begin{itemize}
\raggedright
  \item \emph{Say it:} \emph{a moldura}
  \item \emph{Say it:} \emph{a moldura}, once more
  \item \emph{Say it:} a sentence of your own with \emph{a moldura}
  \item \emph{Recall:} say the Portuguese for canvas, then the Portuguese for easel, then say \emph{a moldura} again
\end{itemize}

\begin{tcolorbox}[breakable,colback=teal!4,colframe=teal!35!black,title={Before you move on}]
What does a moldura mean? (\textquotedblleft{}(picture) frame\textquotedblright{}.) Say it once more.
\end{tcolorbox}
